\documentclass[aspectratio=169,14pt]{beamer}
\usetheme{metropolis}
\usepackage[utf8]{inputenc}
\usepackage[ngerman]{babel}
\usepackage{tikz}
\usetikzlibrary{arrows.meta,positioning,calc,decorations.pathreplacing,shapes,fit,backgrounds}
\usepackage{amsmath,amssymb}
\usepackage{graphicx}
\usepackage{fontawesome5}
\usepackage{xcolor}

\definecolor{accent1}{HTML}{FF6B6B}
\definecolor{accent2}{HTML}{4ECDC4}
\definecolor{accent3}{HTML}{45B7D1}
\definecolor{accent4}{HTML}{96CEB4}
\definecolor{accent5}{HTML}{FFEAA7}
\definecolor{darkbg}{HTML}{2C3E50}
\definecolor{neuronyellow}{HTML}{F39C12}
\definecolor{neuronblue}{HTML}{3498DB}
\definecolor{neurongreen}{HTML}{27AE60}

\setbeamercolor{frametitle}{bg=darkbg,fg=white}
\setbeamercolor{title}{fg=darkbg}
\setbeamercolor{normal text}{fg=darkbg}
\setbeamerfont{frametitle}{size=\large}

\title{Wie funktioniert ChatGPT wirklich?}
\subtitle{Von einfachen Bausteinen zu künstlicher Intelligenz}
\date{}
\author{}

\begin{document}

\begin{frame}[plain]
\maketitle
\end{frame}

% ============================================================
\section{Embeddings -- Recap}
% ============================================================

\begin{frame}{Wörter sind Punkte im Raum}
\begin{columns}
\begin{column}{0.55\textwidth}
Wie gerade gesehen:
\begin{itemize}
\item Jedes Wort hat einen \textbf{Ort} in einem hochdimensionalen Raum
\item Ähnliche Bedeutung = nahe beieinander
\item Diese Punkte sind der \textbf{Startpunkt} für alles Weitere
\end{itemize}
\vspace{0.5cm}
\textbf{Die Frage jetzt:}\\
Was macht die KI mit diesen Punkten?
\end{column}
\begin{column}{0.45\textwidth}
\begin{tikzpicture}[scale=0.85]
\fill[accent3,opacity=0.08] (-0.5,-0.5) rectangle (4.5,4.5);
\node[circle,fill=accent1,inner sep=2.5pt,label={[font=\footnotesize]right:König}] at (3,3.5) {};
\node[circle,fill=accent1,inner sep=2.5pt,label={[font=\footnotesize]right:Königin}] at (3.5,3) {};
\node[circle,fill=accent2,inner sep=2.5pt,label={[font=\footnotesize]left:Hund}] at (0.5,1.5) {};
\node[circle,fill=accent2,inner sep=2.5pt,label={[font=\footnotesize]left:Katze}] at (1,1) {};
\node[circle,fill=neuronyellow,inner sep=2.5pt,label={[font=\footnotesize]below:laufen}] at (2,0.3) {};
\node[circle,fill=neuronyellow,inner sep=2.5pt,label={[font=\footnotesize]above:rennen}] at (2.3,0.8) {};
\draw[->,gray,thin] (0,0) -- (4.5,0);
\draw[->,gray,thin] (0,0) -- (0,4.5);
\end{tikzpicture}
\end{column}
\end{columns}
\end{frame}

% ============================================================
\section{Tokenization}
% ============================================================

\begin{frame}{Wie liest die KI Text?}
\begin{center}
\large Die KI liest weder in Buchstaben noch in Wörtern --\\
sondern in \textbf{Tokens}: häufige Silben und Wortteile.
\end{center}
\vspace{0.7cm}
\begin{center}
\begin{tikzpicture}[node distance=0.15cm]
\node[fill=accent1!30,rounded corners,inner sep=8pt,font=\large] (t1) {un};
\node[fill=accent2!30,rounded corners,inner sep=8pt,font=\large,right=of t1] (t2) {happi};
\node[fill=accent3!30,rounded corners,inner sep=8pt,font=\large,right=of t2] (t3) {ness};
\node[above=0.8cm of t2,font=\Large] {``unhappiness''};
\draw[->,thick] (-0.8,1) -- (t1.north);
\draw[->,thick] (1,1) -- (t2.north);
\draw[->,thick] (2.9,1) -- (t3.north);
\end{tikzpicture}
\end{center}
\vspace{0.7cm}
\begin{columns}
\begin{column}{0.5\textwidth}
\centering
\textcolor{neurongreen}{\faCheck} \textbf{Ca. 50.000--100.000 Tokens}\\[0.2cm]
\small Häufige Wörter = 1 Token\\
Seltene = mehrere Tokens
\end{column}
\begin{column}{0.5\textwidth}
\centering
\textcolor{accent1}{\faExclamationTriangle} \textbf{Konsequenz:}\\[0.2cm]
\small ``Strawberry'' $\rightarrow$ [Str, aw, berry]\\
Die KI ``sieht'' keine Buchstaben!
\end{column}
\end{columns}
\end{frame}

\begin{frame}{Jedes Token wird zum Punkt}
\begin{center}
\begin{tikzpicture}[scale=0.9,>=Stealth]
% Text
\node[font=\large] (text) at (0,3) {``Die Katze sitzt''};
% Tokens
\node[fill=accent1!25,rounded corners,inner sep=6pt] (t1) at (-2.5,1.5) {Die};
\node[fill=accent2!25,rounded corners,inner sep=6pt] (t2) at (0,1.5) {Katze};
\node[fill=accent3!25,rounded corners,inner sep=6pt] (t3) at (2.5,1.5) {sitzt};
\draw[->,thick] (text.south) -- (t1.north);
\draw[->,thick] (text.south) -- (t2.north);
\draw[->,thick] (text.south) -- (t3.north);
% Embeddings
\node[circle,draw,fill=accent1!40,minimum size=0.6cm,font=\tiny] (e1) at (-2.5,-0.5) {};
\node[circle,draw,fill=accent2!40,minimum size=0.6cm,font=\tiny] (e2) at (0,-0.5) {};
\node[circle,draw,fill=accent3!40,minimum size=0.6cm,font=\tiny] (e3) at (2.5,-0.5) {};
\draw[->,thick] (t1.south) -- (e1.north) node[midway,right,font=\tiny] {Embedding};
\draw[->,thick] (t2.south) -- (e2.north);
\draw[->,thick] (t3.south) -- (e3.north);
% Vectors
\node[font=\tiny,below=0.2cm of e1,align=center] {[0.23, $-$0.81,\\1.02, ...]};
\node[font=\tiny,below=0.2cm of e2,align=center] {[0.95, 0.12,\\$-$0.44, ...]};
\node[font=\tiny,below=0.2cm of e3,align=center] {[$-$0.31, 0.67,\\0.88, ...]};
\end{tikzpicture}
\end{center}
\vspace{0.3cm}
\centering
Jedes Token ist jetzt ein Punkt mit \textbf{tausenden Koordinaten}.\\
Diese Punkte werden jetzt \textbf{verarbeitet}.
\end{frame}

% ============================================================
\section{Dense Layer und ReLU}
% ============================================================

\begin{frame}{Baustein 1: Das Dense Layer}
\begin{columns}
\begin{column}{0.5\textwidth}
\textbf{Ein einzelnes Neuron:}
\vspace{0.3cm}
$$y = w_1 \cdot x_1 + w_2 \cdot x_2 + b$$
\vspace{0.3cm}
\begin{tikzpicture}[scale=0.8,>=Stealth]
% Inputs
\node[circle,draw,fill=neuronblue!30,minimum size=0.7cm,font=\tiny] (x1) at (0,1.5) {$x_1$};
\node[circle,draw,fill=neuronblue!30,minimum size=0.7cm,font=\tiny] (x2) at (0,0) {$x_2$};
% Neuron
\node[circle,draw,fill=neuronyellow!50,minimum size=1cm,font=\small] (n) at (3,0.75) {$\Sigma$};
% Output
\node[circle,draw,fill=neurongreen!30,minimum size=0.7cm,font=\tiny] (y) at (5.5,0.75) {$y$};
% Connections
\draw[->,thick] (x1) -- (n) node[midway,above,font=\tiny] {$w_1{=}0.7$};
\draw[->,thick] (x2) -- (n) node[midway,below,font=\tiny] {$w_2{=}{-}0.3$};
\draw[->,thick] (n) -- (y);
% Bias
\node[font=\tiny,below=0.1cm of n] {$b{=}0.1$};
\end{tikzpicture}
\end{column}
\begin{column}{0.5\textwidth}
\textbf{Mehrere Neuronen = ein Layer:}
\vspace{0.3cm}
\begin{tikzpicture}[scale=0.7,>=Stealth]
% Inputs
\foreach \i in {1,2,3} {
\node[circle,draw,fill=neuronblue!30,minimum size=0.6cm,font=\tiny] (x\i) at (0,-\i*1.3+0.5) {$x_{\i}$};
}
% Neurons
\foreach \j in {1,2,3,4} {
\node[circle,draw,fill=neuronyellow!50,minimum size=0.6cm,font=\tiny] (n\j) at (3,-\j*1.0+0.5) {};
}
% Connections
\foreach \i in {1,2,3} {
\foreach \j in {1,2,3,4} {
\draw[->,thin,gray!50] (x\i) -- (n\j);
}
}
% Dots
\node[font=\large] at (3,-4) {$\vdots$};
\node[font=\small,below] at (3,-4.5) {bis zu Millionen!};
\end{tikzpicture}
\end{column}
\end{columns}
\end{frame}

\begin{frame}{Das Problem: Linear kann nur Geraden}
\begin{center}
\begin{tikzpicture}[scale=1.0]
% Left: what linear can do
\begin{scope}[xshift=-4cm]
\draw[->,thin,gray] (-2,0) -- (2.5,0) node[right,font=\tiny] {$x$};
\draw[->,thin,gray] (0,-0.5) -- (0,2.5) node[above,font=\tiny] {$y$};
\draw[very thick,neuronblue] (-1.5,-0.5) -- (2,2);
\node[font=\small\bfseries,below] at (0,-1) {Nur linear: Gerade};
\node[font=\small,neuronblue] at (0,-1.7) {\faCheck\ einfach};
\node[font=\small,accent1] at (0,-2.3) {\faTimes\ langweilig};
\end{scope}
% Right: what we need
\begin{scope}[xshift=4cm]
\draw[->,thin,gray] (-2,0) -- (2.5,0) node[right,font=\tiny] {$x$};
\draw[->,thin,gray] (0,-0.5) -- (0,2.5) node[above,font=\tiny] {$y$};
\draw[very thick,accent1] plot[smooth] coordinates {(-1.5,0.3) (-0.5,2) (0.5,0.5) (1.5,1.8) (2,1)};
\node[font=\small\bfseries,below] at (0,-1) {Realität: Komplex!};
\node[font=\small,accent1] at (0,-1.7) {Wie kriegen wir das hin?};
\end{scope}
% Arrow
\draw[->,very thick,darkbg] (-0.8,1) -- (1.2,1) node[midway,above,font=\small] {???};
\end{tikzpicture}
\end{center}
\end{frame}

\begin{frame}{Die Lösung: ReLU macht Knicke!}
\begin{center}
\begin{tikzpicture}[scale=1.1]
\draw[->,thin,gray] (-3.5,0) -- (3.5,0) node[right,font=\small] {$x$};
\draw[->,thin,gray] (0,-0.8) -- (0,3) node[above,font=\small] {$y$};
% ReLU
\draw[very thick,accent2] (-3.5,0) -- (0,0) -- (3,3);
% Annotations
\node[font=\small,gray] at (-2.5,0.5) {alles $< 0$};
\node[font=\small,gray] at (-2.5,0.1) {wird zu Null};
\draw[->,thick,accent1] (0.15,0.15) -- (1.2,1.5);
\node[font=\small\bfseries,accent1,right] at (1.2,1.5) {Hier ist der Knick!};
\node[font=\normalsize,below] at (0,-1.2) {$\text{ReLU}(x) = \max(0,\; x)$};
\end{tikzpicture}
\end{center}
\vspace{0.3cm}
\centering
\large Simpel? Ja. Aber daraus entsteht alles.
\end{frame}

\begin{frame}{Mehr Neuronen = Mehr Knicke = Jede Kurve}
\begin{center}
\begin{tikzpicture}[scale=0.85]
% 2 neurons
\begin{scope}[xshift=-5cm]
\draw[->,thin,gray] (-1.8,0) -- (1.8,0);
\draw[->,thin,gray] (0,-0.3) -- (0,2.2);
\draw[very thick,accent1] (-1.5,0.5) -- (0,1.8) -- (1.5,0.3);
\fill[accent1] (0,1.8) circle (2.5pt);
\node[font=\small\bfseries,below] at (0,-0.6) {2 Neuronen};
\end{scope}
% 5 neurons
\begin{scope}[xshift=0cm]
\draw[->,thin,gray] (-1.8,0) -- (1.8,0);
\draw[->,thin,gray] (0,-0.3) -- (0,2.2);
\draw[very thick,accent2] (-1.5,0.3) -- (-0.8,1.5) -- (-0.2,0.8) -- (0.5,1.9) -- (1.0,1.2) -- (1.5,1.6);
\foreach \x/\y in {-0.8/1.5, -0.2/0.8, 0.5/1.9, 1.0/1.2} {
\fill[accent2] (\x,\y) circle (2.5pt);
}
\node[font=\small\bfseries,below] at (0,-0.6) {5 Neuronen};
\end{scope}
% Many neurons
\begin{scope}[xshift=5cm]
\draw[->,thin,gray] (-1.8,0) -- (1.8,0);
\draw[->,thin,gray] (0,-0.3) -- (0,2.2);
\draw[very thick,accent3] plot[smooth] coordinates {(-1.5,0.5) (-1,1.2) (-0.5,1.8) (0,1.5) (0.3,0.8) (0.7,1.6) (1.0,2.0) (1.3,1.3) (1.5,1.5)};
\node[font=\small\bfseries,below] at (0,-0.6) {Viele Neuronen};
\node[font=\small,accent3,below] at (0,-1.2) {$\approx$ beliebige Kurve!};
\end{scope}
\end{tikzpicture}
\end{center}
\vspace{0.5cm}
\begin{center}
\tikz{\node[fill=accent4!40,rounded corners,inner sep=10pt,font=\normalsize,align=center]{
\textbf{Universal Approximation Theorem:}\\
Dense + ReLU + Dense kann \emph{jede} Funktion approximieren.\\
Mehr Neuronen = genauer.
};}
\end{center}
\end{frame}

\begin{frame}{Beispiel: Bewässerung und Ernte}
\begin{columns}
\begin{column}{0.55\textwidth}
\begin{tikzpicture}[scale=0.75]
\draw[->,thick] (0,0) -- (7.5,0) node[right,font=\small] {Uhrzeit};
\draw[->,thick] (0,0) -- (0,5) node[above,font=\small] {Ertrag (kg)};
% Time labels
\foreach \x/\l in {1/6h, 2.5/10h, 4/14h, 5.5/18h, 7/22h} {
\node[font=\tiny,below] at (\x,-0.1) {\l};
}
% Learned curve
\draw[very thick,accent2] plot[smooth,tension=0.7] coordinates {(0.5,3) (1,4) (1.8,4.5) (2.5,4.3) (3.2,3.8) (4,2.2) (4.8,3.5) (5.5,4.4) (6.2,3.8) (7,2)};
% Problem zones
\fill[accent1,opacity=0.12] (3.3,0) rectangle (4.7,5);
\fill[gray,opacity=0.1] (6.3,0) rectangle (7.5,5);
\node[font=\tiny,accent1,align=center] at (4,4.7) {Mittag};
\node[font=\tiny,gray,align=center] at (6.9,4.7) {Nacht};
\end{tikzpicture}
\end{column}
\begin{column}{0.45\textwidth}
\textbf{Das Netzwerk lernt:}
\begin{itemize}
\item Mittags wässern $\rightarrow$ schlecht
\item Nachts wässern $\rightarrow$ schlecht
\item Morgens/Abends $\rightarrow$ gut
\end{itemize}
\vspace{0.4cm}
\textbf{Aber es wei\ss{} nicht:}
\begin{itemize}
\item Mittags: Tropfen = Linse $\rightarrow$ Sonnenbrand
\item Nachts: Feuchtigkeit $\rightarrow$ Schnecken
\end{itemize}
\end{column}
\end{columns}
\vspace{0.3cm}
\centering
\small Es findet das Muster -- aber hat \textbf{kein Modell der Welt}.\\
Es ``macht was geht'', ohne zu wissen warum.
\end{frame}

% ============================================================
\section{Der Transformer}
% ============================================================

\begin{frame}{Jetzt: Wie wird daraus eine KI?}
\begin{center}
\large Wir haben jetzt zwei Bausteine:
\end{center}
\vspace{0.5cm}
\begin{columns}
\begin{column}{0.45\textwidth}
\centering
\begin{tikzpicture}
\node[rectangle,draw,fill=accent2!25,rounded corners,minimum width=3.5cm,minimum height=1.5cm,align=center,font=\normalsize] {
\textbf{Dense + ReLU}\\[3pt]
\small Kann alles berechnen
};
\end{tikzpicture}
\end{column}
\begin{column}{0.1\textwidth}
\centering
\Large $+$
\end{column}
\begin{column}{0.45\textwidth}
\centering
\begin{tikzpicture}
\node[rectangle,draw,fill=accent3!25,rounded corners,minimum width=3.5cm,minimum height=1.5cm,align=center,font=\normalsize] {
\textbf{Attention}\\[3pt]
\small Verbindet Wörter
};
\end{tikzpicture}
\end{column}
\end{columns}
\vspace{0.8cm}
\begin{center}
\large $=$ \textbf{Ein Transformer-Layer}
\vspace{0.3cm}

\normalsize Und davon stapelt man 96 Stück übereinander.

\small (GPT-4: ca. 120 Layer, 1.8 Billionen Parameter)
\end{center}
\end{frame}

\begin{frame}{Ein Layer nach dem anderen}
\begin{center}
\begin{tikzpicture}[scale=0.7,>=Stealth,
layer/.style={rectangle,draw,rounded corners,minimum width=8cm,minimum height=0.9cm,font=\small}]
% Input
\node[layer,fill=accent5!40] (input) at (0,0) {Token-Embeddings: Punkte im Raum};
% Layer 1
\node[layer,fill=accent3!20] (a1) at (0,1.5) {Attention: Wörter tauschen Information aus};
\node[layer,fill=accent2!20] (d1) at (0,2.8) {Dense+ReLU: Punkte werden verschoben};
% Layer 2
\node[layer,fill=accent3!20] (a2) at (0,4.3) {Attention};
\node[layer,fill=accent2!20] (d2) at (0,5.6) {Dense+ReLU};
% Dots
\node[font=\Large] at (0,6.7) {$\vdots$};
% Layer N
\node[layer,fill=accent3!20] (aN) at (0,7.8) {Attention};
\node[layer,fill=accent2!20] (dN) at (0,9.1) {Dense+ReLU};
% Output
\node[layer,fill=accent1!30] (out) at (0,10.6) {Vorhersage: nächstes Wort};

\draw[->,thick] (input) -- (a1);
\draw[->,thick] (a1) -- (d1);
\draw[->,thick] (d1) -- (a2);
\draw[->,thick] (a2) -- (d2);
\draw[->,thick] (d2) -- (0,6.3);
\draw[->,thick] (0,7.1) -- (aN);
\draw[->,thick] (aN) -- (dN);
\draw[->,thick] (dN) -- (out);

% Brace
\draw[decorate,decoration={brace,amplitude=8pt},thick] (4.8,1.2) -- (4.8,9.4);
\node[font=\small,align=left,right] at (5.3,5.3) {$\times$ 96\\Layer};

% Annotations
\node[font=\footnotesize,accent3] at (-6,1.5) {Kontext verstehen};
\node[font=\footnotesize,accent2] at (-6,2.8) {Nachdenken};
\end{tikzpicture}
\end{center}
\end{frame}

\begin{frame}{Was macht ein Layer mit den Punkten?}
\begin{center}
\begin{tikzpicture}[scale=0.95]
% Before
\begin{scope}
\node[font=\small\bfseries] at (1.5,3.8) {Vorher};
\fill[accent3,opacity=0.06] (0,0) rectangle (3,3.5);
\fill[accent1] (0.5,2.8) circle (4pt) node[right,font=\tiny] {Die};
\fill[accent2] (2.2,1.5) circle (4pt) node[right,font=\tiny] {Katze};
\fill[accent3] (1.0,0.8) circle (4pt) node[right,font=\tiny] {sitzt};
\end{scope}
% Arrow
\draw[->,very thick,darkbg] (3.5,1.7) -- (5,1.7) node[midway,above,font=\footnotesize] {1 Layer};
% After
\begin{scope}[xshift=5.5cm]
\node[font=\small\bfseries] at (1.5,3.8) {Nachher};
\fill[accent3,opacity=0.06] (0,0) rectangle (3,3.5);
\fill[accent1] (0.8,2.5) circle (4pt) node[right,font=\tiny] {Die};
\fill[accent2] (1.8,2.0) circle (4pt) node[right,font=\tiny] {Katze};
\fill[accent3] (1.3,1.2) circle (4pt) node[right,font=\tiny] {sitzt};
% Show movement
\draw[->,thin,accent2,dashed] (2.2+5.5-5.5,1.5) -- (1.8,2.0);
\end{scope}
% After many layers
\draw[->,very thick,darkbg] (9,1.7) -- (10.5,1.7) node[midway,above,font=\footnotesize] {96 Layer};
\begin{scope}[xshift=11cm]
\node[font=\small\bfseries] at (1.5,3.8) {Am Ende};
\fill[accent3,opacity=0.06] (0,0) rectangle (3,3.5);
\fill[accent1] (1.2,1.0) circle (4pt) node[right,font=\tiny] {Die};
\fill[accent2] (1.5,2.8) circle (4pt) node[right,font=\tiny] {Katze};
\fill[accent3] (2.0,1.8) circle (4pt) node[right,font=\tiny] {sitzt};
\end{scope}
\end{tikzpicture}
\end{center}
\vspace{0.5cm}
\centering
Jeder Layer verschiebt die Punkte ein Stück.\\[0.2cm]
Am Ende kodiert ``Katze'' nicht mehr nur \emph{Katze},\\
sondern \textbf{``die Katze die hier sitzt und über die wir reden''}.
\end{frame}

% ============================================================
\section{Attention im Detail}
% ============================================================

\begin{frame}{Attention: Jedes Wort schaut alle anderen an}
\begin{center}
\begin{tikzpicture}[scale=0.85,>=Stealth]
% Words in a row
\foreach \i/\w/\c in {1/Die/accent1, 2/Katze/accent2, 3/sa\ss/accent3, 4/auf/gray, 5/der/gray, 6/Matte/neuronyellow} {
\node[rectangle,draw,fill=\c!20,rounded corners,minimum width=1.2cm,minimum height=0.7cm,font=\small] (w\i) at (\i*2-2,0) {\w};
}
% Attention from "sa\ss" (w3)
\draw[->,thick,accent3,opacity=0.3] (w3.north) to[bend left=25] (w1.north);
\draw[->,very thick,accent3,opacity=0.9] (w3.north) to[bend left=20] (w2.north);
\draw[->,thick,accent3,opacity=0.2] (w3.north) to[bend right=20] (w4.north);
\draw[->,thick,accent3,opacity=0.15] (w3.north) to[bend right=25] (w5.north);
\draw[->,thick,accent3,opacity=0.7] (w3.north) to[bend right=30] (w6.north);

% Labels
\node[font=\footnotesize,accent3,above=1.5cm of w3,align=center] {``sa\ss'' fragt:\\Wer hat gesessen?\\Wo?};
\node[font=\footnotesize,accent2,above=1.8cm of w2] {\textbf{0.6}};
\node[font=\footnotesize,neuronyellow,above=1.8cm of w6] {\textbf{0.25}};
\end{tikzpicture}
\end{center}
\vspace{0.3cm}
\centering
Jedes Wort stellt eine ``Frage'' an alle anderen.\\
Die Antworten bestimmen, wohin der Punkt verschoben wird.
\end{frame}

\begin{frame}{Attention: Query, Key, Value}
\begin{center}
\begin{tikzpicture}[scale=0.85,>=Stealth]
% Left side: the mechanism
\node[font=\normalsize\bfseries] at (0,4.5) {Jedes Token berechnet drei Dinge:};

\node[rectangle,draw,fill=accent1!20,rounded corners,inner sep=8pt,font=\small,align=center] (q) at (-4,2.5) {\textbf{Query} (Q)\\Was suche ich?};
\node[rectangle,draw,fill=accent2!20,rounded corners,inner sep=8pt,font=\small,align=center] (k) at (0,2.5) {\textbf{Key} (K)\\Was bin ich?};
\node[rectangle,draw,fill=neuronyellow!25,rounded corners,inner sep=8pt,font=\small,align=center] (v) at (4,2.5) {\textbf{Value} (V)\\Mein Inhalt};

% Example
\node[font=\small,align=center] at (0,0.5) {
Query von ``sa\ss'' $\cdot$ Key von ``Katze'' = \textbf{hoher Wert}!\\[4pt]
$\rightarrow$ Also flie\ss t der Value von ``Katze'' in ``sa\ss'' ein.
};

% Formula (small, non-intrusive)
\node[rectangle,draw,fill=gray!10,rounded corners,inner sep=6pt,font=\footnotesize] at (0,-1.2) {
$\text{Attention} = \text{softmax}\!\left(\frac{Q \cdot K^T}{\sqrt{d}}\right) \cdot V$
};
\end{tikzpicture}
\end{center}
\end{frame}

\begin{frame}{Frühe Layer: Tokens zusammenfügen}
\begin{center}
\begin{tikzpicture}[scale=0.9,>=Stealth]
% Before attention (tokenized)
\node[font=\small\bfseries] at (0,3.5) {Tokenisiert:};
\node[fill=accent1!20,rounded corners,inner sep=5pt,font=\small] (t1) at (-3,2.5) {Don};
\node[fill=accent1!20,rounded corners,inner sep=5pt,font=\small] (t2) at (-1.5,2.5) {au};
\node[fill=accent1!20,rounded corners,inner sep=5pt,font=\small] (t3) at (0,2.5) {dam};
\node[fill=accent1!20,rounded corners,inner sep=5pt,font=\small] (t4) at (1.5,2.5) {pf};
\node[fill=accent2!20,rounded corners,inner sep=5pt,font=\small] (t5) at (3,2.5) {schiff};
\node[fill=accent3!20,rounded corners,inner sep=5pt,font=\small] (t6) at (5,2.5) {fahrt};

% After first layers
\draw[->,thick,darkbg] (1,1.8) -- (1,1.2) node[midway,right,font=\footnotesize] {Layer 1--3};

\node[font=\small\bfseries] at (0,0.5) {Nach den ersten Layern:};
\node[fill=accent1!40,rounded corners,inner sep=7pt,font=\small] (g1) at (-1,-.5) {Donaudampf};
\node[fill=accent2!40,rounded corners,inner sep=7pt,font=\small] (g2) at (2.5,-.5) {schifffahrt};

% Attention arrows showing merging
\draw[->,thin,accent1,dashed] (t1.south) -- (g1.north);
\draw[->,thin,accent1,dashed] (t2.south) -- (g1.north);
\draw[->,thin,accent1,dashed] (t3.south) -- (g1.north);
\draw[->,thin,accent1,dashed] (t4.south) -- (g1.north);
\draw[->,thin,accent2,dashed] (t5.south) -- (g2.north);
\draw[->,thin,accent3,dashed] (t6.south) -- (g2.north);

% Explanation
\node[font=\small,gray,align=center] at (1,-2) {Attention in frühen Layern erkennt:\\``Diese Tokens gehören zusammen!''};
\end{tikzpicture}
\end{center}
\end{frame}

\begin{frame}{Spätere Layer: Immer abstrakter}
\begin{center}
\begin{tikzpicture}[scale=0.85,
lbox/.style={rectangle,draw,rounded corners,minimum width=5cm,minimum height=0.9cm,font=\small,align=center}
]
\node[lbox,fill=accent4!20] (l1) at (0,5) {Layer 1--5: Tokens zusammenfügen};
\node[lbox,fill=accent3!20] (l2) at (0,3.5) {Layer 6--20: Grammatik, Satzstruktur};
\node[lbox,fill=accent2!20] (l3) at (0,2) {Layer 21--60: Bedeutung, Fakten};
\node[lbox,fill=accent1!20] (l4) at (0,0.5) {Layer 61--96: Abstrakte Schlussfolgerungen};

\draw[->,thick] (l1) -- (l2);
\draw[->,thick] (l2) -- (l3);
\draw[->,thick] (l3) -- (l4);

% Right side annotations
\node[font=\footnotesize,right=0.5cm,accent4!80!black,align=left] at (l1.east) {``Do-nau'' $\rightarrow$ ``Donau''};
\node[font=\footnotesize,right=0.5cm,accent3!80!black,align=left] at (l2.east) {Subjekt, Verb, Objekt};
\node[font=\footnotesize,right=0.5cm,accent2!80!black,align=left] at (l3.east) {``Hauptstadt von Frankreich''};
\node[font=\footnotesize,right=0.5cm,accent1!80!black,align=left] at (l4.east) {Logik, Planung, Kontext};

% Arrow on left
\draw[->,very thick,gray] (-4,5) -- (-4,0.5);
\node[font=\small,gray,rotate=90] at (-4.7,2.75) {Immer abstrakter};
\end{tikzpicture}
\end{center}
\vspace{0.3cm}
\centering
\small Die späten Layer sind am schwersten zu verstehen --\\
dort passiert das ``Denken'', aber niemand wei\ss{} genau wie.
\end{frame}

\begin{frame}{Nebeneffekt: Warum ``nicht gut'' $\neq$ ``schlecht''}
\begin{center}
\begin{tikzpicture}[scale=0.95]
\fill[accent3,opacity=0.05] (-2.5,-1.5) rectangle (5.5,3.5);
% Points
\node[circle,fill=neurongreen!60,inner sep=4pt,font=\small\bfseries] (gut) at (0,2.5) {gut};
\node[circle,fill=accent1!60,inner sep=4pt,font=\small\bfseries] (schlecht) at (5,0) {schlecht};
\node[circle,fill=neuronyellow!60,inner sep=3pt,font=\footnotesize] (nichtgut) at (0.8,1.8) {};
\node[font=\footnotesize,right] at (0.8,1.5) {``nicht gut''};
% Arrow showing small shift
\draw[->,thick,accent2] (gut) -- (nichtgut);
% Dashed line to where we'd expect it
\draw[->,thin,dashed,gray] (gut) -- (schlecht) node[midway,below,font=\tiny,sloped] {was man erwarten würde};
\end{tikzpicture}
\end{center}
\vspace{0.2cm}
\begin{itemize}
\item Attention bildet \textbf{gewichtete Summen} -- es mittelt zwischen Punkten
\item ``nicht'' verschiebt ``gut'' nur ein \emph{bisschen} weg -- nicht bis zu ``schlecht''
\item Deshalb haben LLMs manchmal Probleme mit Verneinungen
\end{itemize}
\end{frame}

% ============================================================
\section{Next-Token-Prediction}
% ============================================================

\begin{frame}{Am Ende: Welches Wort kommt als nächstes?}
\begin{center}
\begin{tikzpicture}[scale=0.8,>=Stealth]
% Final vector
\node[circle,draw,thick,fill=accent3!30,minimum size=1.3cm,font=\small,align=center] (final) at (0,3) {Finaler\\Punkt};
% Arrow down
\draw[->,very thick] (0,2) -- (0,1);
\node[font=\small,right] at (0.3,1.5) {Vergleich mit \textbf{allen} 100.000 Token-Embeddings};
% Results
\node[font=\small] at (0,0) {Skalarprodukt $\rightarrow$ Ähnlichkeit $\rightarrow$ Wahrscheinlichkeit};
% Bar chart
\draw[fill=neurongreen!60] (-4,-1) rectangle (-3,-3.5);
\draw[fill=neurongreen!40] (-2.5,-1) rectangle (-1.5,-2.8);
\draw[fill=neurongreen!25] (-1,-1) rectangle (0,-1.8);
\draw[fill=neurongreen!15] (0.5,-1) rectangle (1.5,-1.5);
\draw[fill=neurongreen!10] (2,-1) rectangle (3,-1.3);
\node[font=\tiny,below] at (-3.5,-3.5) {Katze};
\node[font=\tiny,below] at (-2,-2.8) {Hund};
\node[font=\tiny,below] at (-0.5,-1.8) {Maus};
\node[font=\tiny,below] at (1,-1.5) {Tisch};
\node[font=\tiny,below] at (2.5,-1.3) {...};
% Percentages
\node[font=\tiny,neurongreen!80!black] at (-3.5,-0.7) {35\%};
\node[font=\tiny,neurongreen!80!black] at (-2,-0.7) {22\%};
\node[font=\tiny,neurongreen!80!black] at (-0.5,-0.7) {10\%};
\end{tikzpicture}
\end{center}
\end{frame}

\begin{frame}{Sampling: Zufall macht kreativ}
\begin{columns}
\begin{column}{0.5\textwidth}
\textbf{Temperature:}
\begin{itemize}
\item Niedrig (0.1): Fast immer das wahrscheinlichste Wort
\item Hoch (1.5): Auch unwahrscheinliche Wörter haben eine Chance
\end{itemize}
\vspace{0.4cm}
\textbf{Top-K Sampling:}
\begin{itemize}
\item Nur die besten K Kandidaten werden betrachtet
\item Rest wird ignoriert
\item Verhindert totalen Unsinn
\end{itemize}
\end{column}
\begin{column}{0.5\textwidth}
\begin{tikzpicture}[scale=0.65]
% Low temperature
\node[font=\small\bfseries] at (1.5,5) {Temp = 0.2};
\draw[fill=accent1!70] (0,0) rectangle (1,4.5);
\draw[fill=accent1!30] (1.3,0) rectangle (2.3,1.2);
\draw[fill=accent1!10] (2.6,0) rectangle (3.6,0.3);
\node[font=\tiny,below] at (0.5,0) {Katze};
\node[font=\tiny,below] at (1.8,0) {Hund};
\node[font=\tiny,below] at (3.1,0) {Tisch};
\end{tikzpicture}
\vspace{0.3cm}
\begin{tikzpicture}[scale=0.65]
% High temperature
\node[font=\small\bfseries] at (1.5,5) {Temp = 1.5};
\draw[fill=accent2!60] (0,0) rectangle (1,2.5);
\draw[fill=accent2!50] (1.3,0) rectangle (2.3,2.0);
\draw[fill=accent2!40] (2.6,0) rectangle (3.6,1.5);
\node[font=\tiny,below] at (0.5,0) {Katze};
\node[font=\tiny,below] at (1.8,0) {Hund};
\node[font=\tiny,below] at (3.1,0) {Tisch};
\end{tikzpicture}
\end{column}
\end{columns}
\end{frame}

% ============================================================
\section{Die Black Box}
% ============================================================

\begin{frame}{Das gro\ss e Geheimnis}
\begin{center}
\begin{tikzpicture}[scale=0.85]
% Black box
\node[rectangle,draw,very thick,fill=darkbg,text=white,minimum width=5cm,minimum height=3cm,rounded corners,align=center] (box) at (0,0) {};
% Question marks inside
\node[font=\Huge,white,opacity=0.8] at (-0.7,0.5) {?};
\node[font=\huge,white,opacity=0.6] at (0.8,0.2) {?};
\node[font=\Large,white,opacity=0.4] at (-0.3,-0.7) {?};
\node[font=\Large,white,opacity=0.5] at (1.2,-0.5) {?};
\node[font=\small,white] at (0,-1.2) {1.800.000.000.000 Parameter};
% Input
\draw[->,very thick,accent2] (-5,0) -- (box.west);
\node[font=\small,accent2,above] at (-5,0.3) {Frage};
% Output
\draw[->,very thick,accent1] (box.east) -- (5,0);
\node[font=\small,accent1,above] at (5,0.3) {Antwort};
\end{tikzpicture}
\end{center}
\vspace{0.5cm}
\begin{center}
\large Wir kennen die Architektur.\\
Wir kennen die Gewichte.\\[0.3cm]
\textbf{Aber wir verstehen nicht, was sie bedeuten.}
\end{center}
\end{frame}

\begin{frame}{Mechanistic Interpretability: Die Box öffnen}
\begin{center}
\large Ein junges Forschungsfeld versucht zu verstehen,\\
\textbf{welche Algorithmen} das Netzwerk intern gelernt hat.
\end{center}
\vspace{0.7cm}
\begin{columns}
\begin{column}{0.5\textwidth}
\centering
\textbf{Die Idee:}\\[0.3cm]
\begin{itemize}
\item Einzelne Neuronen beobachten
\item ``Schaltkreise'' finden
\item Verstehen was \emph{wo} gespeichert ist
\end{itemize}
\end{column}
\begin{column}{0.5\textwidth}
\centering
\textbf{Bisherige Funde:}\\[0.3cm]
\begin{itemize}
\item Induction Heads (Muster-Erkennung)
\item Fakten in bestimmten Layern
\item Mathematische Algorithmen
\item Geographische Karten!
\end{itemize}
\end{column}
\end{columns}
\end{frame}

% ============================================================
\section{Faszinierende Entdeckungen}
% ============================================================

\begin{frame}{Entdeckung 1: Die Weltkarte}
% BILD: Figure 1 aus "Language Models Represent Space and Time"
% Gurnee & Tegmark, 2023 - https://arxiv.org/abs/2310.02207
% Auch: https://twitter.com/wesg52/status/1709590673218875620
\begin{center}
\begin{tikzpicture}[scale=0.65]
\fill[accent3!8] (-7,-4) rectangle (7,4);
\node[font=\large\bfseries] at (0,4.8) {Städte-Embeddings aus einem LLM:};
% Approximate city positions
\fill[accent1] (-4.5,2.5) circle (4pt) node[above right,font=\tiny] {London};
\fill[accent1] (-4,2) circle (4pt) node[below,font=\tiny] {Paris};
\fill[accent1] (-3,2.3) circle (4pt) node[above,font=\tiny] {Berlin};
\fill[accent1] (-3.5,1.3) circle (4pt) node[right,font=\tiny] {Rom};
\fill[accent1] (-3.8,1.8) circle (4pt);
\fill[accent2] (4.5,2) circle (4pt) node[above,font=\tiny] {Tokyo};
\fill[accent2] (3.5,2.3) circle (4pt) node[above,font=\tiny] {Beijing};
\fill[accent2] (3,1.5) circle (4pt) node[right,font=\tiny] {Delhi};
\fill[neurongreen] (-1,-2) circle (4pt) node[below,font=\tiny] {Lagos};
\fill[neurongreen] (0.5,-2.3) circle (4pt) node[below,font=\tiny] {Nairobi};
\fill[neurongreen] (-0.5,-1.5) circle (4pt);
\fill[neuronyellow] (-6,2) circle (4pt) node[left,font=\tiny] {New York};
\fill[neuronyellow] (-6.2,0.8) circle (4pt) node[left,font=\tiny] {Miami};
\fill[neuronyellow] (-5.5,-1) circle (4pt) node[left,font=\tiny] {S\~ao Paulo};
\fill[gray] (4,-3) circle (4pt) node[right,font=\tiny] {Sydney};
\fill[gray] (5,-2.5) circle (4pt) node[right,font=\tiny] {Auckland};
\end{tikzpicture}
\end{center}
\vspace{0.2cm}
\centering
\small Das Modell hat \textbf{nie eine Karte gesehen} -- nur Text!\\
Aber die effizienteste Art, Entfernungen zu kodieren, \emph{ist} eine Karte.\\[0.2cm]
{\tiny Gurnee \& Tegmark (2023): ``Language Models Represent Space and Time''}
\end{frame}

\begin{frame}{Entdeckung 2: Modular Addition}
% Quelle: Neel Nanda et al. (2023) "Progress measures for grokking via mechanistic interpretability"
% https://arxiv.org/abs/2301.05217
\begin{columns}
\begin{column}{0.5\textwidth}
\textbf{Aufgabe:}\\
Trainiere ein kleines Netzwerk auf:
$$(a + b) \mod 113$$
\vspace{0.3cm}

\textbf{Was es intern lernt:}
\begin{enumerate}
\item Zahlen als Punkte auf einem \textbf{Kreis} darstellen
\item Addition = \textbf{Rotation}
\item Nutzt \textbf{Fourier-Transformation}!
\end{enumerate}
\vspace{0.3cm}
\small Niemand hat ihm Fourier beigebracht. Es hat den Algorithmus \textbf{selbst erfunden}.
\end{column}
\begin{column}{0.5\textwidth}
\begin{tikzpicture}[scale=0.9]
% Circle
\draw[thick,accent3] (0,0) circle (2.2cm);
% Points on circle
\foreach \i in {0,8,...,104} {
\fill[accent3,opacity=0.5] ({2.2*cos(\i*360/113)},{2.2*sin(\i*360/113)}) circle (2pt);
}
% Highlight two numbers
\fill[accent1] ({2.2*cos(25*360/113)},{2.2*sin(25*360/113)}) circle (5pt);
\node[accent1,font=\small] at ({2.8*cos(25*360/113)},{2.8*sin(25*360/113)}) {$a$};
\fill[accent2] ({2.2*cos(60*360/113)},{2.2*sin(60*360/113)}) circle (5pt);
\node[accent2,font=\small] at ({2.8*cos(60*360/113)},{2.8*sin(60*360/113)}) {$b$};
% Result
\fill[neuronyellow] ({2.2*cos(85*360/113)},{2.2*sin(85*360/113)}) circle (5pt);
\node[neuronyellow,font=\small] at ({2.8*cos(85*360/113)},{2.8*sin(85*360/113)}) {$a{+}b$};
% Rotation arrow
\draw[->,thick,accent1,dashed] ({1.6*cos(25*360/113)},{1.6*sin(25*360/113)}) arc (25*360/113:85*360/113:1.6);
\node[font=\footnotesize,below] at (0,-2.8) {Zahlen auf dem Kreis};
\end{tikzpicture}
\end{column}
\end{columns}
\vspace{0.2cm}
\centering
{\tiny Neel Nanda et al. (2023): ``Progress measures for grokking via mechanistic interpretability''}
\end{frame}

\begin{frame}{Was wir wissen -- und was nicht}
\begin{center}
\begin{tikzpicture}[scale=0.9]
% Pie chart metaphor
\fill[accent1!30] (0,0) -- (0:2.5) arc (0:340:2.5) -- cycle;
\fill[neurongreen!50] (0,0) -- (340:2.5) arc (340:360:2.5) -- cycle;
\node[font=\small,align=center] at (170:1.5) {Noch nicht\\verstanden};
\node[font=\small,align=center,neurongreen!80!black] at (350:1.8) {};
\draw[->,thick] (350:2.8) -- (350:2.2);
\node[font=\footnotesize,right] at (350:2.9) {$\sim$1--5\% verstanden};
\end{tikzpicture}
\end{center}
\vspace{0.3cm}
\begin{columns}
\begin{column}{0.5\textwidth}
\textbf{Offene Fragen:}
\begin{itemize}
\item Haben LLMs ein ``Weltmodell''?
\item Warum halluzinieren sie?
\item Was passiert in den mittleren Layern?
\end{itemize}
\end{column}
\begin{column}{0.5\textwidth}
\textbf{Theorien:}
\begin{itemize}
\item Kompression des Internets
\item Simulation von ``Autoren''
\item Emergente Weltmodelle
\item Sehr gute Statistik
\end{itemize}
\end{column}
\end{columns}
\end{frame}

% ============================================================
\section{Zusammenfassung}
% ============================================================

\begin{frame}{Der komplette Weg -- auf einen Blick}
\begin{center}
\begin{tikzpicture}[scale=0.65,>=Stealth,
step/.style={rectangle,draw,rounded corners,minimum width=5.5cm,minimum height=0.7cm,font=\small,align=center}]
\node[step,fill=accent5!40] (s1) at (0,7.5) {1. Text wird in Tokens zerlegt};
\node[step,fill=accent4!40] (s2) at (0,6) {2. Jedes Token wird zum Punkt im Raum};
\node[step,fill=accent3!25] (s3) at (0,4.5) {3. Attention: Punkte schauen sich um};
\node[step,fill=accent2!25] (s4) at (0,3) {4. Dense+ReLU: Punkte werden verschoben};
\node[step,fill=gray!15] (s5) at (0,1.5) {5. Schritt 3+4 wird $\sim$96 mal wiederholt};
\node[step,fill=accent1!25] (s6) at (0,0) {6. Finaler Punkt $\rightarrow$ nächstes Wort};
\draw[->,thick] (s1) -- (s2);
\draw[->,thick] (s2) -- (s3);
\draw[->,thick] (s3) -- (s4);
\draw[->,thick] (s4) -- (s5);
\draw[->,thick] (s5) -- (s6);
\end{tikzpicture}
\end{center}
\end{frame}

\begin{frame}{Was wir heute gelernt haben}
\begin{enumerate}
\item \textbf{Tokens:} KI liest in Silben-Stücken, nicht in Wörtern
\item \textbf{Dense + ReLU:} Kann mit genug Neuronen \emph{alles} berechnen
\item \textbf{Attention:} Lässt Wörter miteinander kommunizieren
\item \textbf{Layer:} Frühe = einfach, späte = abstrakt
\item \textbf{Vorhersage:} Am Ende wird das wahrscheinlichste nächste Wort gewählt
\item \textbf{Black Box:} Wir verstehen noch fast nichts vom Innenleben
\item \textbf{Aber:} Faszinierende Strukturen emergieren (Weltkarte, Fourier...)
\end{enumerate}
\vspace{0.5cm}
\centering
\tikz{\node[fill=accent3!15,rounded corners,inner sep=10pt,font=\large,align=center]{
Erstaunlich einfache Bausteine.\\
Erstaunlich komplexes Ergebnis.\\
\textbf{Und noch voller Geheimnisse.}
};}
\end{frame}

\begin{frame}[plain]
\begin{center}
\vspace{2cm}
{\Huge Danke!}\\[1cm]
{\Large Fragen?}\\[1.5cm]
{\small
Quellen zum Vertiefen:\\[0.3cm]
\begin{tabular}{ll}
Weltkarte: & Gurnee \& Tegmark (2023), arXiv:2310.02207\\
Mod 113: & Neel Nanda et al. (2023), arXiv:2301.05217\\
Attention: & Vaswani et al. (2017), ``Attention Is All You Need''\\
Interpretability: & Anthropic Research Blog\\
\end{tabular}
}
\end{center}
\end{frame}

\end{document}
